\section{The discovery: memoized inference in the KV cache}
\label{sec:mechanism}

\begin{figure}[tbp]
\centering
\includegraphics[width=\linewidth]{fig2_mechanism}
\caption{\textbf{Four causal probes for memoized inference.} (a) Refreshing the field's own KV recovers
$\approx 0$ of the decision; recomputing the downstream recovers it fully. (b) Recovery is
suffix-concentrated, accruing only as many post-field tokens are patched. (c) The decision reads almost
entirely from downstream notes, not the field token. (d) Once the value is present, override wording is
redundant and ``re-evaluate'' phrasing hurts---the note is a committed conclusion.}
\label{fig:mech}
\end{figure}

\paragraph{Where the conclusion is written and read.} We first state the account that the probes in this
section establish. A transformer prefills the prompt left-to-right, each position attending to all earlier
ones. At a small number of positions \emph{after} the mutable field---\emph{aggregator} tokens such as
punctuation, newlines, and section breaks, which later tokens route attention through---the model does more
than cache the raw token: it computes the field-conditioned answer (e.g.\ ``status is \texttt{shipped}, so
the action is \emph{deny}'') and writes that conclusion into those positions' cached key/value vectors at
mid-to-late layers. These cached vectors are the \emph{notes}, and they lie \emph{downstream} of the field.
At decode time the decision token does not re-derive the answer from the field; it attends back to the note
positions, and a small set of late-layer attention heads read the stored conclusion into the output logit.
Refreshing the field's \emph{own} KV therefore changes little, since the conclusion was written elsewhere and
the field itself accounts for under $1\%$ of the decision. The four probes below establish each link of this
chain causally: \emph{where} the notes lie, that they encode a \emph{conclusion} rather than a copy of the
field, and that the decision \emph{reads} them.

\paragraph{Setup.} We study a minimal but agent-realistic decision: a context contains a policy rule and a
\emph{mutable field} whose value determines the correct action (e.g.\ ``cancel an order only if its status
is \texttt{pending}''; with \texttt{status=shipped} the correct action flips from \emph{cancel} to
\emph{deny}). After prefilling the full context we compare four cache states at the decision token:
\emph{stale} (old field, old downstream), \emph{field-only} (refresh the field's KV, leave the downstream
stale), \emph{full-downstream} (recompute everything after the field), and \emph{oracle} (a clean prefill
of the new value). We report \emph{decision recovery}: the fraction of the oracle's flip that a cache
state reproduces ($0$ = behaves like stale, $1$ = behaves like oracle). Four independent probes converge
on one account (\Cref{fig:mech}); \Cref{app:examples} gives a complete worked example---the verbatim
prompt, the erratum, and recorded model responses.


\paragraph{(1) Locality: the field's own KV barely matters.} Refreshing only the field's KV recovers
essentially none of the decision---field-only recovery is $-0.028$ on Llama-3.1-8B and near zero across
models---whereas recomputing the downstream recovers it fully ($1.0$) (\Cref{fig:mech}a). The field is
read \emph{indirectly}: its causal contribution to the decision is under $1\%$.

\paragraph{(2) Suffix concentration: the effect lives downstream and late.} Sweeping how many downstream
tokens we patch (ranked by causal effect) shows recovery accrues slowly and saturates only as we include
many tokens after the field, with the causal mass concentrated in mid/late layers
(\Cref{fig:mech}b). The information the decision needs is not in the field but distributed over the tokens
that followed it at prefill (\Cref{fig:mech}c).

\paragraph{(3) Linear decodability.} The field-conditioned conclusion is linearly decodable from those
downstream tokens' residual stream at prefill time---i.e.\ the model has already \emph{computed and
written down} the answer, not merely copied the field.

\paragraph{(4) Knockout and dose-response.} Ablating the high-effect downstream tokens flips the decision
back to stale, and the effect grows with the number/position of memoizing tokens, ruling out a diffuse
explanation: specific aggregator/delimiter positions carry the conclusion. This mirrors the delimiter-token
aggregation reported by \citet{lindsey2025biology} for forward planning; here the stored quantity is a
backward-looking, field-conditioned \emph{conclusion}.

\paragraph{What the note contains.} If the note were a verbatim copy of the field, restating the value
would suffice and emphatic wording would be inert. Instead, a wording ablation (\Cref{fig:mech}d) shows
that once the corrected value is present, the override phrasing is \emph{redundant} (bare value, tagged
update, and explicit override all reach $\approx 1.0$ P(safe)), while aggressive ``disregard your earlier
conclusion and re-evaluate'' phrasing actively \emph{hurts} ($0.81$). The note behaves like a committed
conclusion that a late, salient correction can overwrite---but that confrontational instructions can
destabilize. We name the phenomenon \textbf{attention-mediated memoized inference}. The remainder of this
section stress-tests the account---further controls, more architecture families, and a component-level
circuit---before the rest of the paper builds on it.

\paragraph{Separating the stored conclusion from the field content.} A linear probe can \emph{decode} the
conclusion from the downstream tokens, but decodability alone does not show the decision \emph{uses} it---the
same tokens also encode the field's raw content, and a probe finds both. Four causal controls (three models
each: Qwen3-8B/4B and Llama-3.1-8B; \Cref{app:deepmech}) close that gap; the causal patch, not the probe, is
the instrument that separates the two.

\textbf{(i) It is the conclusion, not the field content.} We hold the field value byte-identical and flip a
single rule token (a polarity \emph{trigger}) so the \emph{conclusion} inverts while the field \emph{content}
stays constant. Transplanting the downstream notes then carries the entire flipped conclusion (recovery
$0.998$--$1.009$), whereas patching the changed rule token carries none ($-0.007$ to $+0.007$). Since the
content never changed, the notes cannot merely be re-encoding it.

\textbf{(ii) The note is written before it is read.} Within the single prefill pass the conclusion becomes
decodable on the downstream aggregator about twelve layers \emph{earlier} (relative depth $0.31$--$0.39$)
than the point at which the decision token commits to its answer (depth $0.73$--$0.77$): the note already
exists at prefill, ahead of the read.

\textbf{(iii) A few specific tokens carry it.} Transplanting the eight \emph{highest-effect} downstream
positions recovers $0.74$--$0.79$ of the decision, while eight \emph{random} downstream positions recover
${\le}0.035$: the conclusion lives in a few specific aggregator tokens, not a diffuse code.

\textbf{(iv) The decision follows the note, even when the note is wrong.} Injecting a \emph{false} note (the
opposite conclusion's downstream KV) into an otherwise-consistent cache makes the decision follow the written
note \emph{against its own live field} (recovery $\approx1.0$); a handful of note tokens suffice. The account
also holds off the synthetic template: field-only recovery stays ${\approx}0$ for multi-hop reasoning and
free-form conversational phrasing, and is bounded only for near-verbatim attribute lookup, where the field is
itself partly a copy (\Cref{app:deepmech}).

\paragraph{Not a model-family artifact: four families.} To rule out the aggregator-token account being
a Qwen3/Llama tokenizer artifact, we replicate five probes (the locality probe and the four deep
controls) on two further architecture families,
\textbf{Gemma-2-9B} and \textbf{Mistral-7B} (a tokenizer-robust readout; for Gemma-2 we keep its
attention/logit soft-capping intact). Every result holds: field-only recovery $0.005$/$0.137$ vs.\
full-downstream $1.0$; the conclusion/content dissociation (trigger-only $\approx0$ with the field held
identical vs.\ notes $\approx1.0$); top-$8$ vs.\ random-$8$ specificity ($0.95$ vs.\ $0.48$); false-note
injection ($0.98$--$1.0$, follow-rate $1.0$); and write-before-read timing (write depth $0.19$--$0.26$ vs.\
decision commit $0.47$--$0.48$). The mechanism is consistent across \emph{four} families (Qwen3, Llama-3.1,
Gemma-2, Mistral; \Cref{app:deepmech}).

\paragraph{A component-level circuit.} The probes above show \emph{where} the conclusion is stored; five
further interventions (\Cref{app:circuit}; replicated across Llama-3.1, Qwen3, Gemma-2, and Mistral) show
\emph{which components} write and read it. The pattern is \emph{distributed write, concentrated read}.

\textbf{Write (many components, mid layers).} Mid-layer \emph{attention}---not the MLPs---does most of the
writing onto the aggregator tokens ($\approx0.6$ of the write on Llama). The conclusion is stored
\emph{redundantly}: it is easy to \emph{decode} (a trained sparse-autoencoder feature separates the two
conclusions at AUC $1.0$), yet its causal content is spread across ${\sim}10$--$30$ features along a shared
low-rank direction (that direction carries $\approx25\times$ the decision-effect of a random one).

\textbf{Read (few components, late layers).} A small, nameable set of late \emph{read heads}---those
attending from the decision token back to the aggregators---funnels the stored conclusion into the output
logit ($12$ such heads recover $0.78$ of the decision; the same number of random heads recover $\approx0$).

\textbf{Control.} Causal scrubbing confirms the note alone governs the decision: resampling it to the
opposite conclusion flips the decision, while resampling everything else leaves it unchanged. As at the
token level, the single feature that best \emph{decodes} the conclusion is not by itself \emph{causally}
sufficient---decodability and causation come apart, now among individual features. With the mechanism
established, the rest of the paper turns it into capabilities: the cache is editable (\Cref{sec:editable})
and composable (\Cref{sec:composable}).
